% Source: https://artofproblemsolving.com/community/c3h1386128_this_is_easy_but_for_some_reason_i_cant_figure_it_out
% Stable identity: aops-topic-1386128-q01
If $A$ has at most two elements, there are no three distinct elements $x,y,z$ in $A$, so the condition holds automatically (it is vacuously true). Take $A=\{\pi\}$. If $\pi\sqrt M=q$ were rational for a positive integer $M$, then $\pi^2=\dfrac{q^2}{M}$ would be rational, which is false ($\pi$ is transcendental, so no power of it is rational). Hence $A=\{\pi\}$ is a counterexample (as is $\{1,\pi\}$, since $\sqrt M$ and $\pi\sqrt M$ cannot both be rational).

So the statement needs the extra hypothesis that $A$ has \emph{at least three} elements. (For example, the only such set containing $0$ and exactly two other elements is $\{-1,0,1\}$: with $A=\{0,a,b\}$, the conditions give $a^2,b^2\in A$, which forces $\{a,b\}=\{1,-1\}$.)
